\chapter{Wnioski i podsumowanie}
\label{chap:wnioski}

Rozdział podsumowuje pracę: ocenia stopień realizacji postawionych celów, omawia najistotniejsze trudności
napotkane w~trakcie wytwarzania systemu wraz z~przyjętymi rozwiązaniami oraz wskazuje kierunki dalszego rozwoju.

\section{Stopień realizacji celów pracy}
\label{sec:stopien-realizacji-celow}

Cel główny pracy - zaprojektowanie, implementacja oraz wdrożenie w~chmurze Amazon Web Services systemu obsługi
zamówień opartego na architekturze mikroserwisowej, został osiągnięty. Powstał działający system, którego
poprawność i~wydajność zweryfikowano w~docelowym środowisku chmurowym (rozdział~\ref{chap:testy-i-analiza}).
Realizację celów szczegółowych, postawionych w~sekcji~\ref{sec:cel-i-zakres-pracy}, podsumowuje
tabela~\ref{tab:realizacja-celow}.

\begin{table}[H]
    \centering
    \caption{Realizacja celów szczegółowych pracy}
    \label{tab:realizacja-celow}
    \begin{tabular}{p{10.5cm}c}
        \toprule
        \textbf{Cel szczegółowy} & \textbf{Realizacja} \\
        \midrule
        Analiza teoretycznych podstaw mikroserwisów i~technologii chmurowych & rozdz.\ \ref{chap:podstawy-mikroserwisow}--\ref{chap:aws} \\
        Projekt modelu domenowego, architektury i~kontraktów komunikacyjnych & rozdz.\ \ref{chap:projekt-systemu} \\
        Implementacja mikroserwisów (zamówienia, produkty, płatności) & rozdz.\ \ref{chap:implementacja-mikroserwisow} \\
        Implementacja aplikacji webowej (portal) & rozdz.\ \ref{chap:implementacja-portalu} \\
        Budowa środowiska chmurowego jako infrastruktury jako kod & rozdz.\ \ref{chap:infrastruktura} \\
        Testy systemu oraz analiza działania i~kosztów & rozdz.\ \ref{chap:testy-i-analiza} \\
        \bottomrule
    \end{tabular}
\end{table}

Zakres pracy zrealizowano zgodnie z~założeniami, z~jednym świadomym ograniczeniem. Spośród pierwotnie rozważanych
usług pomocniczych zaimplementowano trzy serwisy realizujące rdzeń procesu zakupowego, pozostawiając usługi
powiadomień i~analityki jako kierunek dalszego rozwoju (sekcja~\ref{sec:rekomendacje-rozwoju}). Decyzja ta
pozwoliła skupić wysiłek na~kompletności pionowej systemu. Od~interfejsu użytkownika, przez logikę biznesową,
po~infrastrukturę. Zamiast na~zwielokrotnianiu podobnych do~siebie usług.

\section{Napotkane trudności i sposoby ich rozwiązania}
\label{sec:napotkane-trudnosci}

Wytwarzanie systemu wiązało się z~szeregiem problemów, których rozwiązania okazały się pouczające.

\textbf{Limit adresów IP na~węzłach.} Dobór niewielkich instancji \code{t3.small} ze~względów kosztowych
zderzył się z~ograniczeniem wtyczki sieciowej VPC~CNI. Domyślnie węzeł takiej klasy mieści zaledwie~11~podów,
co~uniemożliwiało uruchomienie wszystkich komponentów systemu. Problem rozwiązano włączeniem delegacji prefiksów
(sekcja~\ref{sec:utworzenie-eks}), wielokrotnie podnoszącej ten~limit.

\textbf{Brak komponentu metrics-server.} Automatyczne skalowanie oparte na~zużyciu procesora wymaga komponentu
\code{metrics-server}, który nie~wchodzi w~skład domyślnej instalacji~EKS. Wobec~czego HorizontalPodAutoscaler
pozostawał bezczynny, raportując nieznane wykorzystanie zasobów. Zachowując zasadę utrzymania całości
infrastruktury jako~kod, komponent zainstalowano jako~zarządzany dodatek~EKS, a~nie~ręcznie
(sekcja~\ref{sec:skalowanie-i-cykl-zycia}).

\textbf{Strojenie progu skalowania.} Pierwotny próg HPA na~poziomie~80\% wykorzystania procesora okazał się
zbyt~wysoki. W~teście \emph{stress} skalowanie nie~uruchomiło się mimo~znaczącego obciążenia
(sekcja~\ref{subsec:skalowanie-pod-obciazeniem}). Na~podstawie obserwacji próg obniżono do~60\%, tak~aby
dokładanie replik wyprzedzało nasycenie. Co~zilustrowało wartość empirycznego strojenia parametrów wobec
przyjmowania ich~,,na~wyczucie''.

\textbf{Bezpieczne odizolowanie warstwy aplikacyjnej.} Wystawienie systemu na~świat wyłącznie przez~CloudFront,
przy~jednoczesnym uniemożliwieniu dostępu do~load balancera z~pominięciem warstwy brzegowej, wymagało połączenia
dwóch mechanizmów. Ograniczenia grupy bezpieczeństwa do~zakresów adresowych CloudFront oraz~tajnego nagłówka
weryfikacyjnego wstrzykiwanego przez~CDN i~wymaganego przez~regułę load balancera (sekcja~\ref{sec:routing-ruchu}).

\textbf{Spójność danych w~systemie rozproszonym.} Utrzymanie spójności stanu rozproszonego pomiędzy~usługami
rozwiązano dwojako. Warunkowymi, niepodzielnymi operacjami bazodanowymi odpornymi na~współbieżność (sekcje~\ref{sec:product-service},
\ref{sec:payment-service}). Oraz~idempotentnym przetwarzaniem zarówno żądań~HTTP, jak~i~zdarzeń Kafki.

\section{Rekomendacje dotyczące dalszego rozwoju systemu}
\label{sec:rekomendacje-rozwoju}

Zrealizowany system stanowi kompletną, działającą całość, lecz~jako~projekt o~charakterze demonstracyjnym pozostawia
przestrzeń do~rozwoju w~kierunku wdrożenia produkcyjnego. Najistotniejsze rekomendacje są~następujące.

\begin{itemize}
    \item \textbf{Domknięcie choreografii zdarzeń.} Należy uczynić order-service konsumentem tematu
    \code{product.events}, eliminując przejściowe rozwiązanie, w~którym product-service zapisuje status zamówienia
    bezpośrednio (sekcja~\ref{sec:product-service}). Co~przywróci pełną wyłączność własności danych.

    \item \textbf{Rozdzielenie warstwy danych.} Przydzielenie każdej usłudze własnej bazy oraz~zastąpienie
    transakcyjnej, międzydomenowej finalizacji płatności wzorcem sagi przywróciłoby pełną zgodność ze~wzorcem
    \emph{database-per-service}.

    \item \textbf{Wysoka dostępność komponentów stanowych.} Pojedynczy broker Kafka oraz~jednostrefowa baza danych
    powinny zostać zastąpione konfiguracjami wysokodostępnymi (klaster Kafka lub~Amazon~MSK, tryb~Multi-AZ~RDS),
    a~wewnętrzne proxy zreplikowane i~objęte autoskalowaniem.

    \item \textbf{Rozbudowa funkcjonalna.} Naturalnym rozszerzeniem jest~implementacja usług powiadomień
    i~analityki. Dzięki~architekturze sterowanej zdarzeniami, mogą~dołączyć jako~kolejni konsumenci
    istniejących zdarzeń bez~zmian w~usługach produkujących.

    \item \textbf{Optymalizacje wydajnościowe i~kosztowe.} W~razie~wzrostu obciążenia procesora wskazane jest
    buforowanie wyniku weryfikacji tokenów~JWT (sekcja~\ref{subsec:profilowanie}).
\end{itemize}

\section{Podsumowanie końcowe}
\label{sec:podsumowanie-koncowe}

W~ramach pracy zaprojektowano, zaimplementowano i~wdrożono w~chmurze publicznej kompletny system obsługi zamówień
oparty na~architekturze mikroserwisowej. Powstał system złożony z~trzech mikroserwisów w~języku~Go, komunikujących
się synchronicznie poprzez~REST oraz~asynchronicznie poprzez~zdarzenia domenowe. Powstał z aplikacji webowej~Vue, oraz~brokera komunikatów. 
Wszystkie te komponenty uruchomionych na~klastrze Kubernetes w~usłudze Amazon~EKS. Moduł perzystencji danych został postawiony w usłudze
Amazon RDS. Infrastruktura została zdefiniowana w~całości jako~kod i~objęta monitoringiem.

Praca wykazała, że~zbudowanie tej~klasy systemu rozproszonego skalowalnego, obserwowalnego i~bezpiecznego
jest~dziś w~zasięgu pojedynczego inżyniera. Głównym kosztem architektury mikroserwisowej pozostaje nie~tyle
implementacja samych usług, co~złożoność operacyjna ich~środowiska. Sieci, orkiestracji, wdrażania i~obserwowalności.
To~właśnie zarządzane usługi chmurowe i~podejście infrastruktury jako~kod okazały się~czynnikiem, który~tę~złożoność
czyni możliwą do~opanowania. Przeprowadzone testy potwierdziły, że~zrealizowany system mimo~celowo minimalnych
zasobów, obsługuje setki żądań na~sekundę bezbłędnie. Jego architektura, wraz ze~wskazanymi kierunkami rozwoju,
stanowi solidną podstawę do~ewentualnego wdrożenia produkcyjnego.